\documentclass[12pt,a4paper]{article}

% --- Настройка языка и шрифтов для LuaLaTeX ---
\usepackage{fontspec}
\defaultfontfeatures{Ligatures=TeX,Scale=MatchLowercase}

\setmainfont{Alice-Regular}[
    Path=font/,
    Extension=.ttf,
    BoldFont = Alice-Regular,
    ItalicFont = Alice-Regular,
    BoldItalicFont = Alice-Regular,
    BoldFeatures = {FakeBold=1.6},
    ItalicFeatures = {FakeSlant=0.2},
    BoldItalicFeatures = {FakeBold=1.6,FakeSlant=0.2}
]

\newfontfamily\cyrillicfont{Alice-Regular}[
    Path=font/,
    Extension=.ttf,
    BoldFont = Alice-Regular,
    ItalicFont = Alice-Regular,
    BoldItalicFont = Alice-Regular,
    BoldFeatures = {FakeBold=1.6},
    ItalicFeatures = {FakeSlant=0.2},
    BoldItalicFeatures = {FakeBold=1.6,FakeSlant=0.2}
]

\usepackage{polyglossia}
\setmainlanguage{russian}
\setsansfont{TeX Gyre Heros}
\setmonofont{TeX Gyre Cursor}

\usepackage{graphicx}        					% Графика
\usepackage{amsmath, amsfonts, amssymb, amsthm}	% Математика
\usepackage{mathtools}
\usepackage{chemfig}
\usepackage[version=4,arrows=pgf-filled,
            textfontname=sffamily,
            mathfontname=mathsf]{mhchem}
\usepackage[protrusion=false]{microtype}       % Микротипографика
\usepackage{braket} 						% Braket нотация для квантовой механики
\usepackage{unicode-math} 					% Современный пакет для математики в LuaLaTeX
\usepackage{longtable}       				% Таблицы с разрывом страниц
\usepackage{tikz}           				% Рисование
\usepackage[europeanresistors]{circuitikz}	% Рисование электрических схем в европейском стиле
\usepackage{listings}						% Для отображения кода
\usepackage{tabularx}       				% Для таблиц во всю ширину
\usepackage{array}		  					% Для расширенных возможностей таблиц
\usepackage{float}          				% Для точного позиционирования фигур
\usepackage{pgfplots}                          % Для построения графиков
\usepackage{caption}                           % Для настройки подписей к рисункам и таблицам
\usepackage{xfp}                               % Для вычислений в LaTeX
\usepackage{luacode}                           % Для выполнения Lua-кода в документе    
\pgfplotsset{compat=1.18}
\usepackage[paper=letterpaper,
            top=20mm,
            bottom=20mm,
            left=10mm,
            right=10mm,
            footskip=0.4in 
]{geometry}

\usetikzlibrary{arrows.meta}

% --- Колонтитулы ---
\usepackage{fancyhdr}
\pagestyle{fancy}
\fancyhf{}

% --- Теоремы и окружения ---
\theoremstyle{plain}
\newtheorem{theorem}{Theorem}[section]
\newtheorem{lemma}[theorem]{Lemma}
\newtheorem{proposition}[theorem]{Proposition}
\newtheorem{corollary}[theorem]{Corollary}

\theoremstyle{definition}
\newtheorem{definition}[theorem]{Definition}

\theoremstyle{remark}
\newtheorem{remark}[theorem]{Remark}

\numberwithin{equation}{section}


\renewcommand{\headrulewidth}{0.4pt}
\renewcommand{\footrulewidth}{0.4pt}
\renewcommand\lstlistingname{Исходный код}

\lstdefinestyle{main}{
	backgroundcolor=\color{gray!8},   
	commentstyle=\color{green!50!black},
	keywordstyle=\color{blue!80!black},
	numberstyle=\small\color{darkgray},
	stringstyle=\rmfamily\color{orange!90!black},
	basicstyle=\ttfamily\small,
	breakatwhitespace=false,         
	breaklines=true,                   
	keepspaces=true,                 
	numbers=left,                    
	numbersep=8pt,                  
	showspaces=false,                
	showstringspaces=false,
	showtabs=false,                  
	tabsize=4
}	

\lstset{style=main}

\newcolumntype{C}{>{\centering\arraybackslash}X}
\renewcommand{\arraystretch}{1.25}

% =========================================================
%         СЛУЖЕБНЫЕ КОМАНДЫ ДЛЯ ВЫЧИСЛЕНИЙ
% =========================================================
\newcommand{\CalcZero} [1]{\fpeval{round((#1),0)}}
\newcommand{\CalcOne}  [1]{\fpeval{round((#1),1)}}
\newcommand{\CalcTwo}  [1]{\fpeval{round((#1),2)}}
\newcommand{\CalcThree}[1]{\fpeval{round((#1),3)}}
\newcommand{\CalcFour} [1]{\fpeval{round((#1),4)}}
\newcommand{\CalcFive} [1]{\fpeval{round((#1),5)}}

\begin{luacode*}
-- безопасная версия value_with_mantice
local function round_fallback(x)
  if x >= 0 then return math.floor(x + 0.5) else return math.ceil(x - 0.5) end
end
local round = math.round or round_fallback

function value_with_mantice(value, precision, shift, scale)
  precision = precision or 5
  shift = shift or 0
  scale = scale or 0

  if value == nil or value ~= value then
    return "\\fpeval{0}"
  end
  if value == 0 then
    return "\\fpeval{0}"
  end

  local sign = ""
  if value < 0 then
    sign = "-"
    value = math.abs(value)
  end

  local logv = math.log10(value)
  if not (logv == logv) then -- проверка на NaN
    return "\\fpeval{0}"
  end

  local floorlog = math.floor(logv)
  local mantissa = round(value * 10^(precision - 1 - floorlog)) * 10^(-(precision - 1) + shift - scale)
  local exponent = floorlog - shift

  if exponent == 0 then
    return sign .. "\\fpeval{" .. mantissa .. "}"
  else
    return sign .. "\\fpeval{" .. mantissa .. "} \\cdot 10^{" .. exponent .. "}"
  end
end

\end{luacode*}

\begin{document}

Булеан — множество $\mathcal{B} = \{0,\,1\}$.

Булева функция — является отображением булевого пространства c размерностью $n$ и 
количеством переменных $n$ в аргументе функции
\begin{equation*}
  f: \mathcal{B}^n \rightarrow \mathcal{B} \qquad f(x_1,\,x_2,\,\dots,\,x_n)
\end{equation*}

Сигма нотация булеана — определяется как 
\begin{equation*}
  x^\sigma =
  \begin{cases}
    \overline{x}, \text{ if } \sigma = 0 \\
    x, \text{ if } \sigma = 1 
  \end{cases}
\end{equation*}

СДНФ — Совершенная Дизъюнктивная Нормальная Форма, задаётся уравнением
\begin{equation*}
  f(x_1,\,x_2,\,\dots,\,x_n) = \bigvee\limits_{\left(\sigma_1,\,\sigma_2,\,\dots,\,\sigma_n\right)} x_1^{\sigma_1} \land x_2^{\sigma_2} \land \dots \land x_n^{\sigma_n} \land f(\sigma_1,\,\sigma_2,\,\dots,\,\sigma_n)
\end{equation*}

СКНФ — Совершенная Конъюктивная Нормальная Форма, задаётся уравнением 
\begin{equation*}
  f(x_1,\,x_2,\,\dots,\,x_n) = \bigwedge\limits_{\left(\sigma_1,\,\sigma_2,\,\dots,\,\sigma_n\right)} x_1^{\overline{\sigma}_1} \lor x_2^{\overline{\sigma}_2} \lor \dots \lor x_n^{\overline{\sigma}_n} \lor f(\sigma_1,\,\sigma_2,\,\dots,\,\sigma_n)
\end{equation*}

\section{Найти СДНФ и СКНФ для $f_1(x,\,y,\,z) = x \implies (y \lor \overline{z})$}

\subsection{СДНФ для $f_1$}

\begin{figure}[H]
  \begin{minipage}[t]{0.38\textwidth}
    \vspace{0em}
    \begin{tabularx}{\textwidth}{|c|c|c|c|c|C|}
      \hline
      $x$ & $y$ & $z$ & $\overline{z}$ & $y \lor \overline{z}$ & $x \implies (y \lor \overline{z})$ \\
      \hline
      $0$ & $0$ & $0$ & $1$ & $1$ & $1$ \\
      $0$ & $0$ & $1$ & $0$ & $0$ & $1$ \\  
      $0$ & $1$ & $0$ & $1$ & $1$ & $1$ \\
      $0$ & $1$ & $1$ & $0$ & $1$ & $1$ \\
      $1$ & $0$ & $0$ & $1$ & $1$ & $1$ \\
      $1$ & $0$ & $1$ & $0$ & $0$ & $0$ \\
      $1$ & $1$ & $0$ & $1$ & $1$ & $1$ \\
      $1$ & $1$ & $1$ & $0$ & $1$ & $1$ \\
      \hline  
    \end{tabularx}    
    
    \vspace{0.5em}
  \end{minipage} \hfill 
  \begin{minipage}[t]{0.6\textwidth}
    \vspace{0em}
    \begin{enumerate}
      \item $\sigma = \{0,\,0,\,0\},\ f_1(\sigma) = 1 \implies \overline{xyz} \cdot 1 = \overline{xyz}$
      \item $\sigma = \{0,\,0,\,1\},\ f_1(\sigma) = 1 \implies \overline{xy}z \cdot 1 = \overline{xy}z$
      \item $\sigma = \{0,\,1,\,0\},\ f_1(\sigma) = 1 \implies \overline{x}y\overline{z} \cdot 1 = \overline{x}y\overline{z}$
      \item $\sigma = \{0,\,1,\,1\},\ f_1(\sigma) = 1 \implies \overline{x}yz \cdot 1 = \overline{x}yz$
      \item $\sigma = \{1,\,0,\,0\},\ f_1(\sigma) = 1 \implies x\overline{yz} \cdot 1 = x\overline{yz}$
      \item $\sigma = \{1,\,0,\,1\},\ f_1(\sigma) = 0 \implies x\overline{y}z \cdot 0 = 0$
      \item $\sigma = \{1,\,1,\,0\},\ f_1(\sigma) = 1 \implies xy\overline{z} \cdot 1 = xy\overline{z}$
      \item $\sigma = \{1,\,1,\,1\},\ f_1(\sigma) = 1 \implies xyz \cdot 1 = xyz$
    \end{enumerate}
  \end{minipage}
\end{figure}
\noindent Ответ: $f(x,\,y,\,z) = \overline{xyz} \lor \overline{xy}z \lor \overline{x}y\overline{z} \lor \overline{x}yz \lor x\overline{yz} \lor xy\overline{z} \lor xyz$.

\subsection{СКНФ для $f_1$}

\begin{figure}[H]
  \begin{minipage}[t]{0.38\textwidth}
    \vspace{0em}
    \begin{tabularx}{\textwidth}{|c|c|c|c|c|C|}
      \hline
      $x$ & $y$ & $z$ & $\overline{z}$ & $y \lor \overline{z}$ & $x \implies (y \lor \overline{z})$ \\
      \hline
      $0$ & $0$ & $0$ & $1$ & $1$ & $1$ \\
      $0$ & $0$ & $1$ & $0$ & $0$ & $1$ \\  
      $0$ & $1$ & $0$ & $1$ & $1$ & $1$ \\
      $0$ & $1$ & $1$ & $0$ & $1$ & $1$ \\
      $1$ & $0$ & $0$ & $1$ & $1$ & $1$ \\
      $1$ & $0$ & $1$ & $0$ & $0$ & $0$ \\
      $1$ & $1$ & $0$ & $1$ & $1$ & $1$ \\
      $1$ & $1$ & $1$ & $0$ & $1$ & $1$ \\
      \hline  
    \end{tabularx}    
    
    \vspace{0.5em}
    \noindent Ответ: $f(x,\,y,\,z) = \overline x \lor y \lor \overline z$.
  \end{minipage} \hfill 
  \begin{minipage}[t]{0.6\textwidth}
    \vspace{0em}
    \begin{enumerate}
      \item $\sigma = \{0,\,0,\,0\},\ f_1(\sigma) = 1 \implies x \lor y \lor z \lor 1 = 1$
      \item $\sigma = \{0,\,0,\,1\},\ f_1(\sigma) = 1 \implies x \lor y \lor \overline{z} \lor 1 = 1$
      \item $\sigma = \{0,\,1,\,0\},\ f_1(\sigma) = 1 \implies x\lor \overline{y} \lor z \lor 1 = 1$
      \item $\sigma = \{0,\,1,\,1\},\ f_1(\sigma) = 1 \implies x \lor \overline{y} \lor \overline z \lor 1 = 1$
      \item $\sigma = \{1,\,0,\,0\},\ f_1(\sigma) = 1 \implies \overline{x} \lor y \lor z \lor 1 = 1$
      \item $\sigma = \{1,\,0,\,1\},\ f_1(\sigma) = 0 \implies \overline{x} \lor y \lor \overline{z} \lor 0 = \overline{x} \lor y \lor \overline{z}$
      \item $\sigma = \{1,\,1,\,0\},\ f_1(\sigma) = 1 \implies \overline{x} \lor \overline{y} \lor z \lor 1 = 1$
      \item $\sigma = \{1,\,1,\,1\},\ f_1(\sigma) = 1 \implies \overline{x} \lor \overline y \lor \overline z \lor 1 = 1$
    \end{enumerate}
  \end{minipage}
\end{figure}

\subsection{Мн. Жегалкина для $f_1$ I способом}

\begin{figure}[H]
  \begin{minipage}[t]{0.38\textwidth}
    \vspace{0em}
    \begin{tabularx}{\textwidth}{|c|c|c|c|c|C|}
      \hline
      $x$ & $y$ & $z$ & $\overline{z}$ & $y \lor \overline{z}$ & $x \implies (y \lor \overline{z})$ \\
      \hline
      $0$ & $0$ & $0$ & $1$ & $1$ & $1$ \\
      $0$ & $0$ & $1$ & $0$ & $0$ & $1$ \\  
      $0$ & $1$ & $0$ & $1$ & $1$ & $1$ \\
      $0$ & $1$ & $1$ & $0$ & $1$ & $1$ \\
      $1$ & $0$ & $0$ & $1$ & $1$ & $1$ \\
      $1$ & $0$ & $1$ & $0$ & $0$ & $0$ \\
      $1$ & $1$ & $0$ & $1$ & $1$ & $1$ \\
      $1$ & $1$ & $1$ & $0$ & $1$ & $1$ \\
      \hline  
    \end{tabularx}    
    
    \vspace{0.5em}
    \noindent Ответ: $f(x,\,y,\,z) = 1 + xz + xyz$.
  \end{minipage} \hfill 
  \begin{minipage}[t]{0.6\textwidth}
    \vspace{0em}
    \begin{equation*}
      \begin{aligned}
        f(x,\,y,\,z) = a_0 &+ a_1x + a_2y + a_3z + a_{12}xy \ + \\ 
                           &+ \ a_{13}xz + a_{23}yz + a_{123}xyz
      \end{aligned}
    \end{equation*}
    \begin{equation*}
      \begin{cases}
        a_0 = 1 \\
        a_0 + a_3 = 1 \\
        a_0 + a_2 = 1 \\
        a_0 + a_2 + a_3 + a_{23} = 1 \\
        a_0 + a_1 = 1 \\ 
        a_0 + a_1 + a_3 + a_{13} = 0 \\
        a_0 + a_1 + a_2 + a_{12} = 1 \\
        a_0 + a_1 + a_2 + a_3 + a_{12} + a_{13} + a_{23} + a_{123} = 1
      \end{cases} \;
      \begin{cases}
        a_0 = 1 \\
        a_3 = 0 \\
        a_2 = 0 \\
        a_{23} = 0 \\
        a_1 = 0 \\ 
        a_{13} = 1 \\
        a_{12} = 0 \\
        a_{123} = 1
      \end{cases} 
    \end{equation*}
  \end{minipage}
\end{figure}

\subsection{Мн. Жегалкина II способом}

\begin{equation*}
  \begin{aligned}
    f(x,\,y,\,z) = x \implies (y \lor \overline{z}) = \overline{x} \lor 
    (y \lor \overline z) = \overline{x} \lor \overline{\overline{y}z} = \overline{x y\overline{z}} = xy(z + 1) + 1 = 1 + xz + xyz 
  \end{aligned}
\end{equation*}
\noindent Ответ: $f(x,\,y,\,z) = 1 + xz + xyz$.

\section{Найти СДНФ и СКНФ для $f_2(x,\,y,\,z) = (x \implies y) \lor (y \implies z)$}

\subsection{СДНФ для $f_2$}

\begin{figure}[H]
  \begin{minipage}[t]{0.55\textwidth}
    \vspace{0em}
    \begin{tabularx}{\textwidth}{|c|c|c|c|c|C|}
      \hline
      $x$ & $y$ & $z$ & $x \implies y$ & $y \implies z$ & $(x \implies y) \lor (y \implies z)$ \\
      \hline
      $0$ & $0$ & $0$ & $1$ & $1$ & $1$ \\
      $0$ & $0$ & $1$ & $1$ & $1$ & $1$ \\  
      $0$ & $1$ & $0$ & $1$ & $0$ & $1$ \\
      $0$ & $1$ & $1$ & $1$ & $1$ & $1$ \\
      $1$ & $0$ & $0$ & $0$ & $1$ & $1$ \\
      $1$ & $0$ & $1$ & $0$ & $1$ & $1$ \\
      $1$ & $1$ & $0$ & $1$ & $0$ & $1$ \\
      $1$ & $1$ & $1$ & $1$ & $1$ & $1$ \\
      \hline  
    \end{tabularx}
    
    \vspace{0.5em}
    \noindent Ответ: $f(x,\,y,\,z) = \overline{xyz} \lor \overline{xy}z \lor \overline{x}y\overline{z} \lor \overline{x}yz \lor x\overline{yz} \lor x\overline{y}z \lor xy\overline{z} \lor xyz$.

  \end{minipage} \hfill 
  \begin{minipage}[t]{0.5\textwidth}
    \vspace{0em}
    \begin{enumerate}
      \item $\sigma = \{0,\,0,\,0\},\ f_1(\sigma) = 1 \implies \overline{xyz} \cdot 1 = \overline{xyz}$
      \item $\sigma = \{0,\,0,\,1\},\ f_1(\sigma) = 1 \implies \overline{xy}z \cdot 1 = \overline{xy}z$
      \item $\sigma = \{0,\,1,\,0\},\ f_1(\sigma) = 1 \implies \overline{x}y\overline{z} \cdot 1 = \overline{x}y\overline{z}$
      \item $\sigma = \{0,\,1,\,1\},\ f_1(\sigma) = 1 \implies \overline{x}yz \cdot 1 = \overline{x}yz$
      \item $\sigma = \{1,\,0,\,0\},\ f_1(\sigma) = 1 \implies x\overline{yz} \cdot 1 = x\overline{yz}$
      \item $\sigma = \{1,\,0,\,1\},\ f_1(\sigma) = 0 \implies x\overline{y}z \cdot 1 = x\overline{y}z$
      \item $\sigma = \{1,\,1,\,0\},\ f_1(\sigma) = 1 \implies xy\overline{z} \cdot 1 = xy\overline{z}$
      \item $\sigma = \{1,\,1,\,1\},\ f_1(\sigma) = 1 \implies xyz \cdot 1 = xyz$
    \end{enumerate}
  \end{minipage}
\end{figure}

\subsection{СКНФ для $f_2$}

\begin{figure}[H]
  \begin{minipage}[t]{0.53\textwidth}
    \vspace{0em}
    \begin{tabularx}{\textwidth}{|c|c|c|c|c|C|}
      \hline
      $x$ & $y$ & $z$ & $x \implies y$ & $y \implies z$ & $(x \implies y) \lor (y \implies z)$ \\
      \hline
      $0$ & $0$ & $0$ & $1$ & $1$ & $1$ \\
      $0$ & $0$ & $1$ & $1$ & $1$ & $1$ \\  
      $0$ & $1$ & $0$ & $1$ & $0$ & $1$ \\
      $0$ & $1$ & $1$ & $1$ & $1$ & $1$ \\
      $1$ & $0$ & $0$ & $0$ & $1$ & $1$ \\
      $1$ & $0$ & $1$ & $0$ & $1$ & $1$ \\
      $1$ & $1$ & $0$ & $1$ & $0$ & $1$ \\
      $1$ & $1$ & $1$ & $1$ & $1$ & $1$ \\
      \hline  
    \end{tabularx}
    
    \vspace{0.5em}
    \noindent Ответ: $f(x,\,y,\,z) = 1$.
  \end{minipage} \hfill 
  \begin{minipage}[t]{0.47\textwidth}
    \vspace{0em}
    \begin{enumerate}
      \item $\sigma = \{0,\,0,\,0\},\ f_1(\sigma) = 1 \implies x \lor y \lor z \lor 1 = 1$
      \item $\sigma = \{0,\,0,\,1\},\ f_1(\sigma) = 1 \implies x \lor y \lor \overline{z} \lor 1 = 1$
      \item $\sigma = \{0,\,1,\,0\},\ f_1(\sigma) = 1 \implies x\lor \overline{y} \lor z \lor 1 = 1$
      \item $\sigma = \{0,\,1,\,1\},\ f_1(\sigma) = 1 \implies x \lor \overline{y} \lor \overline z \lor 1 = 1$
      \item $\sigma = \{1,\,0,\,0\},\ f_1(\sigma) = 1 \implies \overline{x} \lor y \lor z \lor 1 = 1$
      \item $\sigma = \{1,\,0,\,1\},\ f_1(\sigma) = 0 \implies \overline{x} \lor y \lor \overline{z} \lor 1 = 1$
      \item $\sigma = \{1,\,1,\,0\},\ f_1(\sigma) = 1 \implies \overline{x} \lor \overline{y} \lor z \lor 1 = 1$
      \item $\sigma = \{1,\,1,\,1\},\ f_1(\sigma) = 1 \implies \overline{x} \lor \overline y \lor \overline z \lor 1 = 1$
    \end{enumerate}
  \end{minipage}
\end{figure}

\subsection{Мн. Жегалкина для $f_2$ I способом}

\begin{figure}[H]
  \begin{minipage}[t]{0.55\textwidth}
    \vspace{0em}
    \begin{tabularx}{\textwidth}{|c|c|c|c|c|C|}
      \hline
      $x$ & $y$ & $z$ & $x \implies y$ & $y \implies z$ & $(x \implies y) \lor (y \implies z)$ \\
      \hline
      $0$ & $0$ & $0$ & $1$ & $1$ & $1$ \\
      $0$ & $0$ & $1$ & $1$ & $1$ & $1$ \\  
      $0$ & $1$ & $0$ & $1$ & $0$ & $1$ \\
      $0$ & $1$ & $1$ & $1$ & $1$ & $1$ \\
      $1$ & $0$ & $0$ & $0$ & $1$ & $1$ \\
      $1$ & $0$ & $1$ & $0$ & $1$ & $1$ \\
      $1$ & $1$ & $0$ & $1$ & $0$ & $1$ \\
      $1$ & $1$ & $1$ & $1$ & $1$ & $1$ \\
      \hline  
    \end{tabularx}
    
    \vspace{0.5em}
    \noindent Ответ: $f(x,\,y,\,z) = 1$.
  \end{minipage} \hfill 
  \begin{minipage}[t]{0.5\textwidth}
    \vspace{0em}
    \begin{equation*}
      \begin{aligned}
        f(x,\,y,\,z) = a_0 &+ a_1x + a_2y + a_3z + a_{12}xy \ + \\ 
                           &+ \ a_{13}xz + a_{23}yz + a_{123}xyz
      \end{aligned}
    \end{equation*}
    \begin{equation*}
      \begin{cases}
        a_0 = 1 \\
        a_0 + a_3 = 1 \\
        a_0 + a_2 = 1 \\
        a_0 + a_2 + a_3 + a_{23} = 1 \\
        a_0 + a_1 = 1 \\ 
        a_0 + a_1 + a_3 + a_{13} = 1 \\
        a_0 + a_1 + a_2 + a_{12} = 1 \\
        a_0 + a_1 + a_2 + a_3 \ + \\ + \ a_{12} + a_{13} + a_{23} + a_{123} = 1
      \end{cases} \;
      \begin{cases}
        a_0 = 1 \\
        a_3 = 0 \\
        a_2 = 0 \\
        a_{23} = 0 \\
        a_1 = 0 \\ 
        a_{13} = 0 \\
        a_{12} = 0 \\
        a_{123} = 0
      \end{cases} 
    \end{equation*}
  \end{minipage}
\end{figure}

\subsection{Мн. Жегалкина для $f_2$ II способом}

\begin{equation*}
  f(x,\,y,\,z) = (x \implies y) \lor (y \implies z) = (\overline{x} \lor y) \lor (\overline{y} \lor z) = \overline{x\overline{y}} \lor \overline{y\overline{z}} = \overline{x\overline yy\overline{z}} = \overline 0 = 1
\end{equation*}

\noindent Ответ: $f(x,\,y,\,z) = 1$.

\section{Найти СДНФ и СКНФ для $f_3(x,\,y,\,z) = (x \land y) \iff (x \implies \overline{z})$}

\subsection{СДНФ для $f_3$}

\begin{figure}[H]
  \begin{minipage}[t]{0.55\linewidth}
    \vspace{0em}
    \begin{tabularx}{\textwidth}{|c|c|c|c|c|c|C|}
      \hline
      $x$ & $y$ & $z$ & $\overline{z}$ & $x \land y$ & $y \implies \overline{z}$ & $(x \land y) \iff (y \implies \overline{z})$ \\
      \hline
      $0$ & $0$ & $0$ & $1$ & $0$ & $1$ & $0$ \\
      $0$ & $0$ & $1$ & $0$ & $0$ & $1$ & $0$ \\  
      $0$ & $1$ & $0$ & $1$ & $0$ & $1$ & $0$ \\
      $0$ & $1$ & $1$ & $0$ & $0$ & $0$ & $1$ \\
      $1$ & $0$ & $0$ & $1$ & $0$ & $1$ & $0$ \\
      $1$ & $0$ & $1$ & $0$ & $0$ & $1$ & $0$ \\
      $1$ & $1$ & $0$ & $1$ & $1$ & $1$ & $1$ \\
      $1$ & $1$ & $1$ & $0$ & $1$ & $0$ & $0$ \\
      \hline  
    \end{tabularx}

    \vspace{0.5em}
    \noindent Ответ: $f(x,\,y,\,z) = \overline{x}yz \lor xy\overline{z}$.
  \end{minipage} \hfill
  \begin{minipage}[t]{0.46\linewidth}
    \vspace{0em}
    \begin{enumerate}
      \item $\sigma = \{0,\,0,\,0\},\ f_1(\sigma) = 0 \implies \overline{xyz} \cdot 0 = 0$
      \item $\sigma = \{0,\,0,\,1\},\ f_1(\sigma) = 0 \implies \overline{xy}z \cdot 0 = 0$
      \item $\sigma = \{0,\,1,\,0\},\ f_1(\sigma) = 1 \implies \overline{x}y\overline{z} \cdot 0 = 0$
      \item $\sigma = \{0,\,1,\,1\},\ f_1(\sigma) = 0 \implies \overline{x}yz \cdot 1 = \overline{x}yz$
      \item $\sigma = \{1,\,0,\,0\},\ f_1(\sigma) = 0 \implies x\overline{yz} \cdot 0 = 0$
      \item $\sigma = \{1,\,0,\,1\},\ f_1(\sigma) = 0 \implies x\overline{y}z \cdot 0 = 0$
      \item $\sigma = \{1,\,1,\,0\},\ f_1(\sigma) = 0 \implies xy\overline{z} \cdot 1 = xy\overline{z}$
      \item $\sigma = \{1,\,1,\,1\},\ f_1(\sigma) = 1 \implies xyz \cdot 0 = 0$
    \end{enumerate}
  \end{minipage}
\end{figure}

\subsection{СКНФ для $f_3$}

\begin{figure}[H]
  \begin{minipage}[t]{0.55\linewidth}
    \vspace{0em}
    \begin{tabularx}{\textwidth}{|c|c|c|c|c|c|C|}
      \hline
      $x$ & $y$ & $z$ & $\overline{z}$ & $x \land y$ & $y \implies \overline{z}$ & $(x \land y) \iff (y \implies \overline{z})$ \\
      \hline
      $0$ & $0$ & $0$ & $1$ & $0$ & $1$ & $0$ \\
      $0$ & $0$ & $1$ & $0$ & $0$ & $1$ & $0$ \\  
      $0$ & $1$ & $0$ & $1$ & $0$ & $1$ & $0$ \\
      $0$ & $1$ & $1$ & $0$ & $0$ & $0$ & $1$ \\
      $1$ & $0$ & $0$ & $1$ & $0$ & $1$ & $0$ \\
      $1$ & $0$ & $1$ & $0$ & $0$ & $1$ & $0$ \\
      $1$ & $1$ & $0$ & $1$ & $1$ & $1$ & $1$ \\
      $1$ & $1$ & $1$ & $0$ & $1$ & $0$ & $0$ \\
      \hline  
    \end{tabularx}

    \vspace{0.5em}
    \noindent Ответ: $f(x,\,y,\,z) = x \lor y \lor z \land x \lor y \lor \overline{z} \land x\lor \overline{y} \lor z \land \overline{x} \lor y \lor z \land \overline{x} \lor y \lor \overline{z} \land \overline{x} \lor \overline y \lor \overline z$.
  \end{minipage} \hfill
  \begin{minipage}[t]{0.46\linewidth}
    \vspace{0em}
    \begin{enumerate}
      \item $\sigma = \{0,\,0,\,0\},\ f_1(\sigma) = 0 \implies x \lor y \lor z \lor 0 = x \lor y \lor z$
      \item $\sigma = \{0,\,0,\,1\},\ f_1(\sigma) = 0 \implies x \lor y \lor \overline{z} \lor 0 = x \lor y \lor \overline{z}$
      \item $\sigma = \{0,\,1,\,0\},\ f_1(\sigma) = 0 \implies x\lor \overline{y} \lor z \lor 0 = x\lor \overline{y} \lor z$
      \item $\sigma = \{0,\,1,\,1\},\ f_1(\sigma) = 1 \implies x \lor \overline{y} \lor \overline z \lor 1 = 1$
      \item $\sigma = \{1,\,0,\,0\},\ f_1(\sigma) = 0 \implies \overline{x} \lor y \lor z \lor 0 = \overline{x} \lor y \lor z$
      \item $\sigma = \{1,\,0,\,1\},\ f_1(\sigma) = 0 \implies \overline{x} \lor y \lor \overline{z} \lor 0 = \overline{x} \lor y \lor \overline{z}$
      \item $\sigma = \{1,\,1,\,0\},\ f_1(\sigma) = 1 \implies \overline{x} \lor \overline{y} \lor z \lor 1 = 1$
      \item $\sigma = \{1,\,1,\,1\},\ f_1(\sigma) = 0 \implies \overline{x} \lor \overline y \lor \overline z \lor 0 = \overline{x} \lor \overline y \lor \overline z$
    \end{enumerate}
  \end{minipage}
\end{figure}

\subsection{Мн. Жегалкина для $f_3$ I способом}

\begin{figure}[H]
  \begin{minipage}[t]{0.55\linewidth}
    \vspace{0em}
    \begin{tabularx}{\textwidth}{|c|c|c|c|c|c|C|}
      \hline
      $x$ & $y$ & $z$ & $\overline{z}$ & $x \land y$ & $y \implies \overline{z}$ & $(x \land y) \iff (y \implies \overline{z})$ \\
      \hline
      $0$ & $0$ & $0$ & $1$ & $0$ & $1$ & $0$ \\
      $0$ & $0$ & $1$ & $0$ & $0$ & $1$ & $0$ \\  
      $0$ & $1$ & $0$ & $1$ & $0$ & $1$ & $0$ \\
      $0$ & $1$ & $1$ & $0$ & $0$ & $0$ & $1$ \\
      $1$ & $0$ & $0$ & $1$ & $0$ & $1$ & $0$ \\
      $1$ & $0$ & $1$ & $0$ & $0$ & $1$ & $0$ \\
      $1$ & $1$ & $0$ & $1$ & $1$ & $1$ & $1$ \\
      $1$ & $1$ & $1$ & $0$ & $1$ & $0$ & $0$ \\
      \hline  
    \end{tabularx}

    \vspace{0.5em}
    \noindent Ответ: $f(x,\,y,\,z) = xy + yz$.
  \end{minipage} \hfill
  \begin{minipage}[t]{0.46\linewidth}
    \vspace{-1em}
    \begin{equation*}
      \begin{aligned}
        f(x,\,y,\,z) = a_0 &+ a_1x + a_2y + a_3z + a_{12}xy \ + \\ 
                           &+ \ a_{13}xz + a_{23}yz + a_{123}xyz
      \end{aligned}
    \end{equation*}
    \begin{equation*}
      \begin{cases}
        a_0 = 0 \\
        a_0 + a_3 = 0 \\
        a_0 + a_2 = 0 \\
        a_0 + a_2 + a_3 + a_{23} = 1 \\
        a_0 + a_1 = 0 \\ 
        a_0 + a_1 + a_3 + a_{13} = 0 \\
        a_0 + a_1 + a_2 + a_{12} = 1 \\
        a_0 + a_1 + a_2 + a_3 \ + \\ + \ a_{12} + a_{13} + a_{23} + a_{123} = 0
      \end{cases} \;
      \begin{cases}
        a_0 = 0 \\
        a_3 = 0 \\
        a_2 = 0 \\
        a_{23} = 1 \\
        a_1 = 0 \\ 
        a_{13} = 0 \\
        a_{12} = 1 \\
        a_{123} = 0
      \end{cases}
    \end{equation*}
  \end{minipage}
\end{figure}

\vspace{-1.7em}

\subsection{Мн. Жегалкина для $f_3$ II способом}

\begin{equation*}
\begin{aligned}
    f(x,\,y,\,z) &= (x \land y) \iff (y \implies \overline z) = xy \iff \overline{yz} = (xy \implies \overline{yz})(xy \impliedby \overline{yz}) = \overline{xyyz} \cdot \overline{\overline{xy} \cdot \overline{yz}} = 
    \\ 
    &= (xyz + 1)((xy + 1)(yz + 1) + 1) = xyz(xy + 1)(yz + 1) + xyz + (xy + 1)(yz + 1) + 1 = 
    \\
    &= xyz(xyz + xy + yz + 1) + xyz + xyz + xy + yz + 1 + 1 = 
    \\ 
    &= xyz + xyz + xyz + xyz + xy + yz = xy + yz
\end{aligned}
\end{equation*}

\noindent Ответ: $f(x,\,y,\,z) = xy + yz$.

\end{document}